\documentclass[]{article}

\usepackage{amsmath}
\usepackage{multirow}
\usepackage{natbib}
\usepackage{graphicx} 


\begin{document}

Globular clusters (GCs) are dense, ancient stellar systems that serve as unique fossils of the early Universe. Within the standard cosmological model, their formation is thought to be intimately linked to the most intense star-forming episodes during the assembly of the first galaxies. A fundamental, yet unresolved, question in galaxy formation is what physical mechanism set the characteristic mass and time scales for the birth of these objects. A leading theoretical framework posits that the formation of the most ancient GCs was regulated by the atomic cooling threshold. This model predicts that efficient star formation, and by extension the assembly of GCs, could only occur within dark matter halos massive enough for their gas to cool via atomic hydrogen transitions. This process becomes effective when the halo's virial temperature, $T_{vir}$, exceeds approximately $10^4$ K, thereby establishing a minimum mass for the host environments of the first GCs. This physical threshold, in turn, predicts a characteristic cosmic history for GC formation, providing a falsifiable prediction that has been historically difficult to test at the relevant epochs.

The advent of powerful new observatories has recently made it possible to discover and characterize candidate GCs at the high redshifts where the predictions of the atomic cooling model are most prominent. Two such discoveries offer an unprecedented opportunity to test this framework across a vast expanse of cosmic time: the GEMS system, a collection of 19 proto-GC candidates observed at a redshift of $z=9.625$, and the Sparkler system, a group of 5 GCs observed at $z=1.4$. These two populations, separated by over eight billion years of cosmic evolution, provide distinct snapshots of the GC population. The GEMS clusters probe the era of cosmic dawn, when the first GCs were forming, while the Sparkler clusters represent a more mature population observed closer to the peak of cosmic star formation. If the atomic cooling threshold is a universal regulator, then the inferred formation histories of these two disparate systems should be consistent with a single, underlying theoretical prediction.

In this paper, we test the universality of the atomic cooling threshold by mapping the empirical formation epochs of the GEMS and Sparkler GCs onto the cosmic GC formation rate predicted by the model. For each cluster, we use its published age and observed redshift to calculate its formation redshift, $z_{form}$. This procedure allows us to place both the high-redshift GEMS population and the lower-redshift Sparkler population onto a common theoretical timeline. We find that the formation epochs of the Sparkler clusters align remarkably well with the predicted peak of GC formation activity. In contrast, the GEMS clusters populate the high-redshift tail of the same distribution, a result consistent with an observational selection effect favoring the discovery of the earliest-forming objects at high redshift. Furthermore, our analysis of the cluster metallicities reveals that the GEMS clusters are unexpectedly metal-rich for their early formation, implying they were born within a massive and rapidly enriching host environment. The consistency of these two distinct populations with a single theoretical framework provides strong evidence that the atomic cooling threshold is a primary physical process governing the formation of globular clusters.

\end{document}
                